\chapter{Métodos Analíticos y Numéricos}
\label{appendix:metodos_numericos}

En este apéndice se presentan herramientas matemáticas complementarias utilizadas en el desarrollo de este trabajo. En particular, se incluye una introducción a la aproximación estocástica y un resultado sobre la divergencia de Kullback-Leibler para distribuciones normales multivariadas.

\section{Aproximación estocástica}
La aproximación estocástica se refiere a una amplia clase de métodos estocásticos iterativos que tienen el objetivo de encontrar raíces o resolver problemas de optimización. No obstante, la definición de estas tareas en muchas ocasiones implica estimar funciones desconocidas. \cite{Shiyu2022}

Para motivar el estudio de esta técnica y encontrar su relación con el aprendizaje por refuerzo, se retoma el ejemplo presentado en \cite{Shiyu2022}.

\subsection{Motivación}
Sea $Z$ una variable aleatoria. Se desea obtener una estimación del valor esperado $\mathbb{E}[Z]$. Con el fin de realizar la estimación, se considera un conjunto de muestras i.i.d $\{z_{i}\}$. Entonces, el valor esperado de $Z$ se aproxima por la siguiente expresión,
\begin{align}
    \mathbb{E}[Z] \approx \frac{1}{n}\sum_{i=1}^{n}z_{i}.
    \label{eq:MC_mean}
\end{align}
Esta aproximación corresponde a la estimación Monte Carlo. Por la ley de los grandes números, se sabe que conforme $n$ crece, la aproximación se asemeja al valor esperado real.

El cálculo de la \eqreftext{eq:MC_mean}, tiene una desventaja significativa, ya que antes de la hacer la estimación se requiere tener un número significativo de muestras. Esto representa un inconveniente si el tiempo de muestreo es tardado. Para resolver esta cuestión, se puede calcular la \eqreftext{eq:MC_mean} de forma incremental e iterativa. Esto se hace realizando el siguiente cálculo en cada iteración $k$,
\begin{align}
    w_{k} = \frac{1}{k}\sum_{i=1}^{k}
    z_i \quad \forall k \in \{2, 3, \cdots\}.
\end{align}
La aproximación $w_{k}$ se puede expresar en términos de $w_{k-1}$, como se muestra debajo,
\begin{equation}
    \begin{aligned}
        w_{k} &= \frac{1}{k}\left(z_{k} + \sum_{i=1}^{k-1}z_{i} \right)\\
        &= \frac{1}{k}(z_{k}+(k-1)w_{k-1}) \\
        &= w_{k-1}-\frac{1}{k}(w_{k-1} -z_{k})
    \end{aligned}
    \label{eq:SA_mean}
\end{equation}
La expresión \eqreftext{eq:SA_mean} permite obtener una estimación del promedio desde el primer momento en que se tenga una muestra. Por lo que puede ser utilizada para otra tarea mientras continua la estimación del promedio.

Adicionalmente, se puede considerar una modificación de la regla de optimización en la \eqreftext{eq:SA_mean}, dada por la siguiente expresión,
\begin{align}
    w_{k} &= w_{k-1}-\alpha_{k}(w_{k-1} -z_{k}),
    \label{eq:SA_mean2}
\end{align}
en este caso el término $\alpha_{k}$ remplaza a $1/k$. La sucesión dada por la \eqreftext{eq:SA_mean2} converge bajo condiciones especificas que debe satisfacer la sucesión $\{\alpha_{k}\}$. Este es un caso particular del algoritmo llamado Robbins-Monro.

\subsection{Método Robbins-Monro}
\label{sec:RM}
El algoritmo \acrfull{rm} es un método pionero en el campo de la aproximación estocástica. Tiene aplicaciones en varias áreas como optimización, aprendizaje automático, en particular, en el aprendizaje por refuerzo. Este método permite estimar un parámetro desconocido, mientras que lo actualiza de forma iterativa, utilizando observaciones ruidosas. \cite{Shiyu2022} 

En general, este método busca resolver problemas con la siguiente forma, 
\begin{align}
    g(w) =0,
    \label{eq:RM_problem}
\end{align}
donde $w\in \mathbb{R}$ es la variable para la que se resuelve el problema y $g:\mathbb{R}\rightarrow{\mathbb{R}}$ es una función. Concretamente, el método \acrshort{rm} se especializa en resolver el caso en que la función $g$ es desconocida y no puede ser observada con precisión. En lugar de ello, se asume que las observaciones disponibles tienen un ruido aditivo, como se muestra en la siguiente expresión,
\begin{align}
    \tilde{g}(w, \eta) = g(w) + \eta,
\end{align}
donde $\eta \in \mathbb{R}$ es un error de observación.

Bajo las consideraciones anteriores el método \acrshort{rm} resuelve \eqreftext{eq:RM_problem}, utilizando la siguiente regla de actualización,
\begin{align}
    w_{k+1} &= w_{k} -\alpha_k\tilde{g}(w_k, \eta_k),
    \label{eq:RM_rule}
\end{align}
donde $w_k$ es la k-ésima estimación de la raíz, $\tilde{g}(w_k, \eta_k)$ es su correspondiente observación y $\alpha_k$ es un coeficiente positivo.

El teorema de Robbins-Monro garantiza la convergencia de la sucesión $w_k$ dada por la \eqreftext{eq:RM_rule} bajo ciertas condiciones. En primer lugar establece que $g$ debe ser monótona creciente o no decreciente. En segundo lugar la sucesión $\{\alpha_k\}$ debe satisfacer las siguientes hipótesis,
\begin{align}
    1) \sum_{k=1}^{\infty}\alpha_{k}^2 < \infty \qquad 2)\sum_{k=1}^{\infty} \alpha_{k} \rightarrow \infty.
\end{align}
Según se expone en \cite{Shiyu2022}, la condición $1)$ garantiza que la diferencia $w_{k+1}-w_{k} \rightarrow 0$, es decir que haya convergencia. Mientras que la condición $2)$ asegura que la convergencia no sea tan rápida y lo haga a partir de cualquier valor inicial.

\subsubsection{Aplicación a la estimación del promedio}
Retomando la \eqreftext{eq:SA_mean2}, que es la aplicación del método \acrshort{rm} para la estimación del promedio. Se puede considerar que 
\begin{align}
    g(w) = w - \mathbb{E}[Z].
\end{align}
Al resolver la ecuación $g(w)=0$, se logra encontrar $\mathbb{E}[Z]$. En este caso la observación disponible está dada por 
\begin{align}
    \tilde{g}(w, z) = w-z.
\end{align}
Ya que $z$ es una muestra de $Z$, se obtiene que,
\begin{equation}
    \begin{aligned}
    \tilde{g}(w, \eta) &= w-z \\
    &= w - z + \mathbb{E}[Z] - \mathbb{E}[Z] \\
    &= \underbrace{(w - \mathbb{E}[Z])}_{g(w)} + \underbrace{(\mathbb{E}[Z]-z)}_{\eta}.
    \end{aligned}
\end{equation}



\section{Divergencia de Kullback-Leibler para distribuciones normales multivariadas}

El siguiente teorema y su correspondiente demostración están basados en el resultado presentado en \textit{Kullback-Leibler divergence for the multivariate normal distribution} \cite{soch2023}.

\begin{theorem}
    Sea $\mathbf{x} \in \mathbb{R}^{n}$ un vector aleatorio. Se asume que dos funciones de densidad de distribuciones normales multivariadas $p$ y $q$ describen la probabilidad de distribución de $\mathbf{x}$ de la siguiente manera,
    \begin{align}
        p: & \quad \mathbf{x} \sim \mathcal{N}(\bm{\mu}_1, \bm{\Sigma}_1), \\
        q: & \quad \mathbf{x} \sim \mathcal{N}(\bm{\mu}_2, \bm{\Sigma}_2).
    \end{align}
    Entonces la divergencia de Kullback-Leibler de $p$ con respecto a $q$ está dada por la siguiente expresión,
    \begin{align}
    D_{\text{KL}}(p||q) &= \frac{1}{2}\left[ 
    (\bm{\mu}_1 - \bm{\mu}_2)^{\top}\bm{\Sigma}_2^{-1}(\bm{\mu}_2-\bm{\mu}_1)
    +\text{tr} \left(\bm{\Sigma}_2^{-1}\bm{\Sigma}_1\right)
    - \log|\bm{\Sigma}_1 \bm{\Sigma}_2^{-1}|-n\right].
    \label{KL_normal}
    \end{align}
\end{theorem}

\begin{proof}
    La divergencia KL para dos variables aleatorias continuas
    está dada por,
    \begin{align}
        D_{\text{KL}}(p||q) = \int_{\mathbb{R}^{n}}p(\mathbf{x})\log \frac{p(\mathbf{x})}{q(\mathbf{x})}d\mathbf{x},
    \end{align}
    aplicando esta expresión a las funciones de densidad de las normales multivariables $p$ y $q$, se obtiene que,
\begin{equation} 
    \begin{aligned}
        D_{\text{KL}}(p||q) &= \int_{\mathbb{R}^{n}} 
        \mathcal{N}(\mathbf{x}; \bm{\mu}_1, \bm{\Sigma}_1)
        \log\frac{
        \mathcal{N}(\mathbf{x}; \bm{\mu}_1, \bm{\Sigma}_1)
        }{
        \mathcal{N}(\mathbf{x}; \bm{\mu}_2, \bm{\Sigma}_2)
        }d\mathbf{x} \\
        &= \left\langle 
        \log\frac{
        \mathcal{N}(\mathbf{x}; \bm{\mu}_1, \bm{\Sigma}_1)
        }{
        \mathcal{N}(\mathbf{x}; \bm{\mu}_2, \bm{\Sigma}_2)
        }
        \right\rangle_{p} \\
        &= \left\langle 
        \log \frac{
        \left((2\pi)^{n}|\bm{\Sigma}_1|\right)^{-\frac{1}{2}}
        \cdot 
        \exp\left[- \frac{1}{2}(\mathbf{x} -\bm{\mu}_1)^{\top}\bm{\Sigma}_1^{-1}(\mathbf{x} - \bm{\mu}_1)\right]
        }{
        \left((2\pi)^{n}|\bm{\Sigma}_2|\right)^{-\frac{1}{2}}
        \cdot 
        \exp\left[- \frac{1}{2}(\mathbf{x} -\bm{\mu}_2)^{\top}\bm{\Sigma}_2^{-1}(\mathbf{x} - \bm{\mu}_2)\right]
        }\right\rangle_{p} \\
        &= \frac{1}{2}\left\langle 
        \log\frac{|\bm{\Sigma}_2|}{|\bm{\Sigma}_1|}
        -(\mathbf{x}-\bm{\mu}_1)^{\top}\bm{\Sigma}_{1}^{-1}(\mathbf{x}- \bm{\mu}_1)
        + (\mathbf{x}-\bm{\mu}_2)^{\top}\bm{\Sigma}_{2}^{-1}(\mathbf{x}- \bm{\mu}_2)
        \right\rangle_p.
    \end{aligned}
\end{equation}
Usando el hecho de que para $a \in \mathbb{R}$, $a=\text{tr}(a)$ y aplicando la propiedad $\text{tr}(\mathbf{ABC})=\text{tr}(\mathbf{BCA})$, se tiene que,
\begin{equation}
    \begin{aligned}
        &= \frac{1}{2}\left\langle 
        \log|\bm{\Sigma}_2\bm{\Sigma}_1^{-1}| 
        - \text{tr}\left(\bm{\Sigma}_1^{-1}(\mathbf{x}-\bm{\mu}_1)(\mathbf{x}-\bm{\mu}_1)^{\top} \right)
        + \text{tr}\left(\bm{\Sigma}_2^{-1}(\mathbf{x}-\bm{\mu}_2)(\mathbf{x}-\bm{\mu}_2)^{\top}\right) 
        \right\rangle_p \\
        &= \frac{1}{2}\left\langle 
        \log|\bm{\Sigma}_2\bm{\Sigma}_1^{-1}| 
        - \text{tr}\left(\bm{\Sigma}_1^{-1}(\mathbf{x}-\bm{\mu}_1)(\mathbf{x}-\bm{\mu}_1)^{\top} \right)
        + \text{tr}\left(
        \bm{\Sigma}_2^{-1}(\mathbf{x}\mathbf{x}^{\top} -2 \bm{\mu}_2\mathbf{x}^{\top} + \bm{\mu}_2\bm{\mu}_2^{\top})
        \right)\right\rangle_p.
    \end{aligned}
\end{equation}
    Debido a que la función $\text{tr}$ y el operador $\langle\cdot\rangle_p$ son lineales, se puede realizar el siguiente cambio sobre el lado derecho de la igualdad, 
    \begin{align}
        &= \frac{1}{2}\left(\log|\bm{\Sigma}_2\bm{\Sigma}_1^{-1}|
        - \text{tr}\left(\bm{\Sigma}_{1}^{-1}\left\langle 
        (\mathbf{x}- \bm{\mu}_1)(\mathbf{x}-\bm{\mu}_1)^{\top}
        \right\rangle_{p}\right) 
        + \text{tr}\left(
        \bm{\Sigma}_2^{-1}\left[\left\langle \mathbf{x}\mathbf{x}^{\top}\right\rangle_p 
        - 2\left\langle \bm{\mu}_2\mathbf{x}^{\top}\right\rangle_p + 
        \left\langle \bm{\mu}_2\bm{\mu}_2^{\top}\right\rangle_p\right]
        \right)
        \right)
    \end{align}
    Usando el hecho de que al calcular valor esperado para $p$ de una forma lineal se obtiene la ecuación debajo, 
    \begin{equation}
    \langle \mathbf{A}\mathbf{x}\rangle_p= \mathbf{A}\bm{\mu}_1,
    \end{equation}
    y que el valor esperado para $p$ de una forma cuadrática da lo siguiente,
    \begin{equation}
    \langle \mathbf{x}^{\top}\mathbf{A}\mathbf{x} \rangle_p = \bm{\mu}^{\top}_1\mathbf{A}\bm{\mu}_1 + \text{tr}(\mathbf{A}\bm{\Sigma}_1),
    \end{equation}
    se obtiene la siguiente expresión,
    \begin{equation}
    \begin{aligned}
        2 \cdot D_{\text{KL}}(p||q) &= 
        \log |\bm{\Sigma}_2\bm{\Sigma}_1^{-1}| 
        - \text{tr}\left(\bm{\Sigma}_{1}^{-1}\bm{\Sigma}_1\right) + 
        \text{tr}\left(\bm{\Sigma}_2^{-1}\left[\bm{\Sigma}_1 
        + \bm{\mu}_1\bm{\mu}_1^{\top}
        -2 \bm{\mu}_2\bm{\mu}_1^{\top}
        +\bm{\mu}_2\bm{\mu}_2^{\top}\right]\right) \\
        &= \log |\bm{\Sigma}_2\bm{\Sigma}_1^{-1}|  - \text{tr}(\mathbf{I}_n) 
        + \text{tr}\left( \bm{\Sigma}_2^{-1}\bm{\Sigma}_1\right)  
        +\text{tr}\left(\bm{\Sigma}_2^{-1}\left[
        \bm{\mu}_1\bm{\mu}_1^{\top}
        -2\bm{\mu}_2\bm{\mu}_1^{\top}
        + \bm{\mu}_2\bm{\mu}_2^{\top}\right] \right)\\
        &= \log |\bm{\Sigma}_2\bm{\Sigma}_1^{-1}| - n + 
        \text{tr}\left( \bm{\Sigma}_2^{-1}\bm{\Sigma}_1\right) +
        \text{tr}\left(
        \bm{\mu}_1^{\top}\bm{\Sigma}_2^{-1}\bm{\mu}_1
        -2\bm{\mu}_1^{\top}\bm{\Sigma}_2^{-1}\bm{\mu}_2 
        + \bm{\mu}_2^{\top}\bm{\Sigma}_2^{-1}\bm{\mu}_2
        \right).
    \end{aligned}
    \end{equation}
    Finalmente, factorizando y ordenando los términos se obtiene,
    \begin{align}
    D_{\text{KL}}(p||q)= \frac{1}{2}\left[ 
    (\bm{\mu}_1 - \bm{\mu}_2)^{\top}\bm{\Sigma}_2^{-1}(\bm{\mu}_2-\bm{\mu}_1)
    +\text{tr}\left(\bm{\Sigma}_2^{-1}\bm{\Sigma}_1\right)
    - \log|\bm{\Sigma}_1 \bm{\Sigma}_2^{-1}|
    -n\right].
    \end{align}
\end{proof}



    %\item Differential entropy of the multivariate normal distribution:
%\url{https://statproofbook.github.io/P/mvn-dent.html}
%\item KL divergence between two multivariate Gaussians: 
%\url{https://stats.stackexchange.com/questions/60680/kl-divergence-between-two-multivariate-gaussians}
%\url{https://statproofbook.github.io/P/mvn-kl.html}
%\cite{brunton2019}



